% ACC 2027 / PaperPlaza official class. ieeeconf.cls MUST sit in the project
% root next to this file - Overleaf does not search subfolders for .cls.
\documentclass[letterpaper, 10 pt, conference]{ieeeconf}

\IEEEoverridecommandlockouts   % required before \thanks in ieeeconf
\overrideIEEEmargins           % required to meet the ACC printer margins

\usepackage{amsmath,amssymb,bm}
\let\proof\relax
\let\endproof\relax
\usepackage{amsthm} % for theorem/proof environments
\newtheorem{theorem}{Theorem}
% for theorem/proof environments
\usepackage{booktabs,array,cite}
\usepackage{graphicx}
\usepackage[hidelinks]{hyperref}
% [v9] Certificate constants are generated from the simulation code, not typed:
% experiments/theory_certificate_2026-09-12/compute_local_certificate.py
\input{certificate_macros.tex}
% [v9] Public code link: set once here when the repository is live.
\newcommand{\coderepo}{\url{https://github.com/[to-be-released]}}
% Figures live under figures/..., so keep that on the graphics path.
\graphicspath{{figures/}{figures/disturbance/}{../experiments/nengo_nef_replication_2026-08-29/results/plots/}}

% [v6] TikZ support for the architecture figures.
\usepackage{tikz}
\usepackage{verbatim}
\usetikzlibrary{arrows.meta,positioning,fit,backgrounds,calc}
% Architecture drawing helpers live in architectures/tikz/.
\input{architectures/tikz/arch_style.tex}

\newcommand{\R}{\mathbb R}
\newcommand{\T}{^{\mathsf T}}
\newcommand{\sat}{\operatorname{sat}}
\newcommand{\clip}{\operatorname{clip}}
\newcommand{\sgnp}{\operatorname{sgn}_{+}}
\newtheorem{proposition}{Proposition}
\newtheorem{remark}{Remark}

% [v7] Title now names the mechanism and the learning, and claims the
% controlled study rather than architectural novelty.
\title{\LARGE \bf Online PES Decoder Learning in Neuromorphic Residual Control of UAV Quaternion Tracking}

% ACC/PaperPlaza template author format (ieeeconf.cls does not support IEEEtran blocks).
\author{Gandhi Contreras$^{1}$, Steven Sandoval$^{1}$, and Luis Rodolfo Garcia Carrillo$^{2}$%
\thanks{$^{1}$New Mexico State University (NMSU), Las Cruces, NM, USA.}%
\thanks{$^{2}$ Space Vehicles Directorate, Air Force Research Laboratory (AFRL), Kirtland Air Force Base (KAFB), Albuquerque, NM 87108.
    Email: {\tt\small luis\_rodolfo.garcia\_carrillo@afrl.af.mil}}%
}

\begin{document}
% IEEE reference practice: lists of more than six authors print as "et al."
\bstctlcite{IEEEexample:BSTcontrol}
\maketitle

% NOTE: The official ACC/PaperPlaza template often suppresses page numbers for
% submission. If you need page numbers for internal review, use `plain`.
\pagestyle{plain}
\thispagestyle{plain}

% [v7] Rewritten: leads with the mechanism and the result; scope limits kept
% but stated once rather than after every claim.
%We present a compact neuromorphic residual controller for quaternion attitude tracking with bounded actuator commands. A fixed quaternion PD baseline is augmented by six paired coupled membrane--gate (CMG) units: continuous, non-spiking, conductance-based neurons whose gate-derived activities form a nonlinear basis. Eighteen decoder weights adapt online via prescribed error sensitivity (PES), and the bounded residual is \emph{added} to the saturated baseline so the baseline is recovered exactly at zero authority. We give bounds on internal states, command deviation, and projected decoder weights, and we isolate three commonly conflated mechanisms: intrinsic conductance modulation, decoder learning, and residual authority.
%In ten matched 45-s MuJoCo cases spanning five missions on Skydio X2 and Crazyflie~2.0 models (the latter completed with an explicit four-rotor interface), CMG--PES reduces mean quaternion-angle error by $17.3$--$60.6\%$ relative to PD while respecting actuator limits (nonzero bound occupancy only in brief Hardstep intervals). Direct and Nengo implementations agree to $2.013\times10^{-6}$ relative mean error. In the tested gain regime, most of the improvement is carried by online decoder adaptation; toggling conductance modulation is numerically indistinguishable. A benchmark against LIF, ALIF, Izhikevich, and LIFRate bases shows the same hierarchy and exposes an effort--accuracy tradeoff: families that match CMG--PES on error spend several times its residual authority to do so. Noise reruns keep the reduction range at $18.1$--$55.8\%$ with seed standard deviation below $0.031\%$, and an intermittent-disturbance study shows learned corrections reused across episodes. Results are in simulation and establish boundedness, not closed-loop stability.

\begin{abstract}
Online neuromorphic residual control can augment a classical attitude controller under abrupt missions or unforeseen disturbances without changing its feedback gains or actuator interface. We present six continuous coupled membrane-gate (CMG) units and eighteen decoder weights adapted by prescribed error sensitivity (PES), adding a bounded correction to the saturated quaternion proportional-derivative (PD) torque request before rotor-force allocation. Across ten matched 45-second MuJoCo cases on the Skydio X2 and a Crazyflie~2 derived model, mean quaternion-angle error decreases by 17.3--60.6\%, including approximately 53--57\% reductions on the Full3 and Hard45 multi-axis tracking missions. Direct and Nengo implementations agree closely. Benchmarks against LIF, ALIF, Izhikevich, and LIFRate neuron models expose an effort-accuracy tradeoff: CMG--PES attains better tracking with substantially lower realized residual effort. Under synthetic measurement noise, the error reductions persist, and learned corrections are reused across repeated disturbances. Overall, PD provides a nominal, interpretable reference, while the online residual adds bounded, revocable adaptation that improves robustness to disturbances and noisy measurements.
\end{abstract}

\section{Introduction}

Learning-based augmentation can retain classical feedback while adapting a correction that is difficult to specify \emph{a priori}. We study a bounded torque residual for quaternion attitude tracking of an unmanned aerial vehicle (UAV) with individual rotor-force saturation. The question is whether a small stateful basis improves a fixed PD controller, and whether any improvement comes from its neuron dynamics, decoder learning, or correction magnitude. Bounded authority makes the learning layer switchable; it does not by itself ensure stability or safety.

Prior work includes residual robot learning \cite{johannink2019}, geometric $\mathcal L_1$ quadrotor augmentation with performance guarantees \cite{wu2025l1quad}, and proportional--integral--derivative (PID) attitude control augmented with a proximal policy optimization (PPO) residual \cite{li2027hpappo}.
%
Brain-inspired robust UAS controllers likewise pair model-based feedback with learned terms, from limbic-system designs \cite{rubioscola2023limbic,escareno2024limbic} to neural weight reset under impulsive disturbances \cite{rubioscola2023weightreset}, and perturbation estimation targets the same disturbances on small quadrotors \cite{rubioscola2021mambo}.
Neuromorphic control offers a complementary implementation approach: Nengo supports the Neural Engineering Framework (NEF) \cite{eliasmith2004neuraleng} and online decoder learning \cite{bekolay2014nengo}, most often through the prescribed error sensitivity (PES) rule \cite{macneil2011pes}, an established local, online error-driven update \cite{voelker2015pes}. PES-based adaptive control has been demonstrated on manipulators \cite{dewolf2016arm,dewolf2020nengo}, and in self-tuning spiking PID controllers for free-flying UAS trajectory tracking \cite{olivares2025ifac,olivares2026selftuning}.
Building on that line, our focus is a continuous membrane-gate basis for a bounded quaternion-torque residual, with a fixed classical controller and adaptation confined to its decoder.

We augment a quaternion PD baseline with a compact neuromorphic residual architecture based on coupled conductance-based membrane-gate (CMG) units and PES decoder learning. The paper's contributions are:
\begin{enumerate}
    \item a code-aligned CMG--PES architecture with paired currents, explicit encoding, state-update, and decoder/readout equations, and a fully specified actuator interface (Fig.~\ref{fig:arch});
    \item proofs of internal-state, command-deviation, and weight boundedness, with exact recovery of the PD baseline at zero residual authority, and an allocation-aware local input-to-state certificate whose constants are computed for both vehicles;
    \item a controlled ablation that separates conductance modulation, decoder learning, and output authority, and identifies which mechanism is primarily responsible for the tracking-error reduction; and
    \item robustness evidence that the reduction survives realistic measurement noise, together with a hold-equilibrium study in which the learned correction is reused across repeated disturbance episodes rather than relearned.
\end{enumerate}

  %\item a matched five-mission, two-model study in which the direct and Nengo realizations agree to within $2.013\times10^{-6}$ in the maximum relative mean quaternion-angle tracking error difference across the ten missions;

The CMG unit couples a leaky membrane to a gating state that modulates its effective conductance, linking LIF-like membrane integration with Hodgkin--Huxley inspired gating dynamics \cite{hodgkin1952quantitative}. Although related operating regimes can support spike-like responses, the present controller deliberately exploits the continuous, non-spiking membrane-gate trajectory, which provides a compact set of stateful temporal features for the residual.
The present implementation is digital; the same temporal-feature mechanism holds regardless of computing platform, and no neuromorphic-hardware energy advantage is claimed.

Section~\ref{sec:problem} defines the problem and actuator interface; Section~\ref{sec:architecture} presents the CMG--PES architecture, its boundedness properties, and a local stability certificate; Section~\ref{sec:results} reports the simulation study, ablations, benchmarks, and noise and disturbance experiments; and Section~\ref{sec:conclusion} discusses limitations and future work. Code, model files, and run archives will be available at \coderepo.

% [v6] Architecture overview (TikZ source under architectures/tikz/).

\begin{figure*}
    \centering
    % Trim bottom whitespace in the standalone PDF so the caption sits closer.
    \includegraphics[width=\textwidth,trim=0 10mm 0 0,clip]{architectures/architecture_standalone.pdf}
    \caption{Online PES decoder learning CMG-PES Neuromorphic Residual Control architecture}
    \label{fig:arch}
\end{figure*}

\begin{comment}
\begin{figure*}[t]

\centering
\resizebox{\textwidth}{!}{\input{architectures/tikz/archA2p.tex}}
\caption{CMG--PES residual architecture. The light blocks are the fixed
classical path: error \eqref{eq:theta}, PD law \eqref{eq:pd}, and
allocation \eqref{eq:mixer} are never adapted. The heavy blocks encode the
clipped baseline as paired currents \eqref{eq:encoding}, evolve $2m$
membrane--gate units \eqref{eq:cmg}, and decode a bounded residual
\eqref{eq:readout} that is \emph{added} to the saturated baseline
\eqref{eq:residual}. Only $W_k$ adapts, through \eqref{eq:pes}; $\beta=0$
recovers the baseline exactly.}
\label{fig:arch}
\end{figure*}
\end{comment}

% [v6] fig:scale and fig:models were merged: both showed the same two
% aircraft. The old fig:scale call sites now point here.
% [v10] Shrunk from \figure* to single-column for the 8-page target.
% TODO(Gandhi): the source image still has 3 baked-in panels (CF2, X2,
% relative-scale) sized for full \textwidth; re-render as a 2-panel
% (CF2, X2) image at single-column width before submission, or this will
% be too small to read. Kept full-width content/caption otherwise.
\begin{figure}[t]
\centering
\includegraphics[width=\linewidth]{figures/01_mujoco_missions/fig2_drones_only.pdf}
\caption{Vehicles and simulation models. Left: CF2 hardware and MuJoCo
model. Centre: X2 hardware and MuJoCo model. Right: both geometries at
relative scale. The axis triads are body-frame axes: they rotate with the
vehicle relative to the inertial frame (cf.\ $Q\!:\,\text{body}\to\text{inertial}$,
Sec.~\ref{sec:problem}).}
\label{fig:models}
\end{figure}


%%%%%%%%%%%%%%%%%%%%%%%%%%%%%%%%%%%%%%%%%%%%%%%%%%%%%%%%%%%%%%%%%%%%%%%%%%
% [v6] Nomenclature. Every symbol that appears in an equation is listed
% once here; the bracketed tag is the equation that defines it.

\begin{comment}
\begin{table*}[t]
\caption{Nomenclature. Parenthesised tags give the defining equation;
vectors are bold; $\clip_{[a,b]}$ and $\sat_{\bm\ell}$ act elementwise.}
\label{tab:nomen}\centering\scriptsize
\setlength{\tabcolsep}{4pt}
\renewcommand{\arraystretch}{0.84}
\begin{tabular}{@{}l>{\raggedright\arraybackslash}p{0.185\textwidth}l>{\raggedright\arraybackslash}p{0.185\textwidth}l>{\raggedright\arraybackslash}p{0.185\textwidth}@{}}
\toprule
\multicolumn{6}{@{}l}{\textit{Attitude, plant, and notation}}\\
$q,\;q_d$ & state / reference quaternion & $q_e$ & error quaternion, $q_d\otimes q^{-1}$ & $q_{e,0},\bm q_{e,v}$ & scalar and vector parts\\
$\bm e_q$ & shortest-path error \eqref{eq:theta} & $\theta_q$ & error angle \eqref{eq:theta} & $\bm\omega$ & body angular rate\\
$J$ & inertia tensor & $Q$ & body-to-inertial rotation & $\otimes$ & quaternion product\\
$\bm\omega_d,\bm\omega_e$ & reference and error rate \eqref{eq:errkin} & $R(\cdot)$ & rotation matrix of a quaternion & $z,z_d,\dot z$ & height, reference, rate \eqref{eq:collective}\\
$\odot,\oslash$ & elementwise $\times$, $\div$ & $\sgnp$ & sign, with $\sgnp(0)=1$ & $\sigma$ & logistic function\\
$I_m$ & identity of order $m$ & $\R$ & real numbers & $\Pi$ & Euclidean projection\\
\midrule
\multicolumn{6}{@{}l}{\textit{Fixed baseline and actuator interface}}\\
$\bm\tau_0=\bm u_0$ & baseline torque \eqref{eq:pd} & $K_q,K_\omega$ & PD gains \eqref{eq:pd} & $\bm\tau_{\max}$ & per-axis torque bound\\
$\omega_n,\zeta$ & design frequency, damping & $\bm\ell$ & per-channel command bound & $\sat_{\bm\ell}$ & saturation at $\pm\bm\ell$\\
$B$ & allocation matrix \eqref{eq:allocationmatrix} & $\bm f$ & rotor forces \eqref{eq:mixer} & $T$ & collective thrust \eqref{eq:collective}\\
$\kappa_i$ & rotor reaction coefficient & $m$ & torque channels, $m=3$ & $M$ & vehicle mass\\
$\bm f_{\max}$ & per-rotor force bound & $g$ & gravity, $9.81$\,m/s$^2$ & $\bm 1$ & vector of ones\\
\midrule
\multicolumn{6}{@{}l}{\textit{Coupled membrane--gate unit}}\\
$v_j$ & membrane state \eqref{eq:cmg} & $r_j$ & gate state \eqref{eq:cmg} & $c_j$ & shunt conductance \eqref{eq:cmg}\\
$\bar{\bm u}$ & normalized baseline \eqref{eq:encoding} & $\bm I^{\pm},\bm I$ & paired currents \eqref{eq:encoding} & $I_b,I_g$ & current offset and gain\\
$\tau,\tau_r$ & membrane, gate time constants & $v_{\rm th}$ & gate threshold & $v_{\rm rest},v_{\rm reset}$ & rest level, reversal level\\
$\gamma$ & gate self-coupling & $\alpha_h$ & logistic sharpness & $\chi$ & conductance switch, $\chi\in\{0,1\}$\\
$g_0,g_1$ & conductance gains (distinct from $g$) & $G$ & activity gain, $G=10^5$ & $h$ & integration step\\
$b_{j,k}$ & gate target \eqref{eq:gate_step} & $r^+_{j,k},v^+_{j,k}$ & updated states \eqref{eq:mem_step} & $\lambda_{j,k},d_{j,k}$ & frozen membrane coefficients\\
\midrule
\multicolumn{6}{@{}l}{\textit{Decoder, online learning, and residual authority}}\\
$\bm a_k$ & clipped activity \eqref{eq:readout} & $W_k$ & decoder \eqref{eq:readout} & $W_0$ & initial decoder, $[I_m,-I_m]$\\
$\bm s_k$ & clipped readout \eqref{eq:readout} & $\bm y^*_k$ & PES target \eqref{eq:pes} & $\bm\varepsilon_k$ & PES error \eqref{eq:pes}\\
$k_e,k^*_\omega$ & target gains \eqref{eq:pes} & $\eta$ & per-step learning rate & $\lambda_W$ & per-step decoder decay\\
$\Pi_{\mathcal W},W_{\max}$ & projection and weight bound & $\beta$ & residual authority \eqref{eq:residual} & $\bm u_k,\bm\delta_k$ & applied command, deviation \eqref{eq:residual}\\
\midrule
\multicolumn{6}{@{}l}{\textit{Reported metrics}}\\
$\bar\theta_q$ & mean quaternion-angle error & $S$ & actuator-bound occupancy [\%] & $n$ & runs per reported cell\\
\bottomrule
\end{tabular}
\end{table*}
\end{comment}

\section{Attitude Baseline and Actuator Interface}
\label{sec:problem}
%%%%%%%%%%%%%%%%%%%%%%%%%%%%%%%%%%%%%%%%%%%%%%%%%%%%%%%%%%%%%%%%%%%%%%%%

The objective is to reduce shortest-path attitude error for a prescribed
$q_d(t)$ while respecting rotor-force limits. Baseline and augmented
controllers share the state measurements, height loop, and allocator;
Section~\ref{sec:architecture} modifies only the requested attitude torque.

%\input{figure1_interface.tex}

\subsection{Quaternion Convention and Error}
Let $Q\in\mathbb S^3$ denote the scalar-first Hamilton quaternion representing attitude (body to inertial), and let $\bm\omega$ be the body angular velocity. For a rigid body with inertia $\bm J\succ0$ about the center of mass, expressed in body axes (cf.~Fig.~\ref{fig:models}), the free-flight rotational dynamics are \cite{schaub2018amss}
\begin{align}
 \dot Q&=\tfrac12 Q\otimes[0,\bm\omega\T]\T,\\
 \bm J\dot{\bm\omega}&=\bm\tau_{\rm act}+\bm\tau_{\rm p}
       +\bm\tau_{\rm d}-\bm\omega\times(\bm J\bm\omega).
 \label{eq:plant}
\end{align}
Here $\bm\tau_{\rm p}$ denotes modeled passive torque and $\bm\tau_{\rm d}$ any applied disturbance; the nominal ideal rigid-body model sets both to zero. Contact, if present, contributes an additional external wrench.

We work with the corresponding passive quaternion $q=Q^{-1}$ \cite{diebel2006quaternion}, which satisfies
\begin{equation}
 \dot q=-\tfrac12[0,\bm\omega\T]\T\otimes q.
 \label{eq:passive}
\end{equation}
The desired attitude $q_d$ uses the same convention. Define the attitude error
\begin{align}
 q_e&=q_d\otimes q^{-1}=[q_{e,0},\bm q_{e,v}\T]\T,\\
 \bm e_q&=\sgnp(q_{e,0})\bm q_{e,v},\qquad \sgnp(0)=1,\\
 \theta_q&=2\operatorname{atan2}(\|\bm q_{e,v}\|_2,|q_{e,0}|)
             \in[0,\pi].
 \label{eq:theta}
\end{align}
The reported error $\theta_q$ is the shortest rotation angle, invariant under $q_e\mapsto-q_e$. Sign selection introduces a discontinuity at $q_{e,0}=0$; no continuous feedback globally stabilizes rigid-body attitude \cite{bhat2000topological,chaturvedi2011rigid}, robust global tracking requires hybrid \cite{mayhew2011hybrid} or geometric-algebra \cite{rubioscola2026ga,escamilla2023ga} formulations, and the analysis here is local.

For a desired-frame rate $\bm\omega_d$ satisfying the corresponding passive kinematics, differentiation gives
\begin{align}
 \dot q_e&=\tfrac12q_e\otimes[0,\bm\omega_e\T]\T,\\
 \bm\omega_e&=\bm\omega-R(q_e)\T\bm\omega_d,\qquad
 \dot q_{e,0}=-\tfrac12\bm q_{e,v}\T\bm\omega_e.
 \label{eq:errkin}
\end{align}
These are kinematic identities (no convergence implied). Reference Euler triples parameterize $q_d$ in this passive convention.

\subsection{Physical Inertia and Implemented PD Gains}
When principal moments $\bm j_{\rm p}$ of $\bm J$ are specified in a rotated principal-axis frame, the physical body-frame tensor is
\begin{equation}
 \bm J=R_{BI}\operatorname{diag}(\bm j_{\rm p})R_{BI}\T,
 \label{eq:inertiaframe}
\end{equation}
for the rotation from principal axes to body axes $R_{BI}$. The implementation sets axiswise gains using the principal moments in their reported ordering, $\bm j_{\rm tune}$, not necessarily $\operatorname{diag}(\bm J)$ (Table~\ref{tab:models}).

The fixed quaternion PD baseline is
\begin{align}
 \bm k_q&=2\bm j_{\rm tune}\odot\bm\omega_n^2,\qquad
 \bm k_\omega=2\zeta\bm j_{\rm tune}\odot\bm\omega_n,\\
 \bm u_0&=\sat_{\bm\ell}
       (-\bm k_q\odot\bm e_q-\bm k_\omega\odot\bm\omega),
 \label{eq:pd}
\end{align}
where $\odot$ denotes componentwise multiplication,
$\bm\omega_n=(4,4,3.5)\T$ rad/s, $\zeta=0.95$, and
$\sat_{\bm\ell}$ clips component $i$ to $[-\ell_i,\ell_i]$.
%
The requested-torque box is $\mathcal U=\prod_{i=1}^3[-\ell_i,\ell_i]$;
$\clip_{[a,b]}$ also acts componentwise, and $\oslash$ denotes componentwise division.
We impose componentwise torque limits $\bm\ell\in\mathbb R_{>0}^3$ through $\sat_{\bm\ell}(\cdot)$. In our implementation, $\bm\ell$ is chosen proportional to the diagonal gain-scheduling vector $\bm j_{\rm tune}$ via
$\ell_i=\max(j_{{\rm tune},i}\nu_i,10^{-6}\ \mathrm{N\,m})$,
with acceleration-scale vector $\bm\nu=(12,12,10)\T$ rad/s$^2$.
These controller-side bounds are not the rotor-feasible torque set.
%
The law is quaternion feedback \cite{fresk2013quaternion}; $\bm\omega_n$ is a tuning parameter rather than a claimed set of coupled closed-loop natural frequencies.
Desired-rate feedforward and gyroscopic cancellation are omitted.
The resulting gains are vehicle-specific through $\bm j_{\rm tune}$, even though
$\bm\omega_n$ and $\zeta$ are shared; the same fixed gains are then retained
for PD and its CMG augmentation.
The local result in Section~\ref{sec:local} uses these anisotropic gains; it does not replace them with a scalar-gain controller.
%
\subsection{Requested and Realized Wrench}
For rotor positions $\bm p_i$ in body coordinates, parallel $+z$ thrust
axes, and signed reaction coefficients $\kappa_i$, define
\begin{equation}
 B=\begin{bmatrix}
 1&1&1&1\\p_{1,y}&p_{2,y}&p_{3,y}&p_{4,y}\\
 -p_{1,x}&-p_{2,x}&-p_{3,x}&-p_{4,x}\\
 \kappa_1&\kappa_2&\kappa_3&\kappa_4
 \end{bmatrix}.
 \label{eq:allocationmatrix}
\end{equation}
The signs follow $\bm p_i\times[0,0,f_i]\T$ and the actuator's local reaction torque $\kappa_i f_i$ \cite{garciacarrillo2013quad}. Here $B$ maps rotor thrusts to total thrust and body torque; $\kappa_i$ has units of meters, not angular velocity. For four rotors $B^{-1}$ equals the Moore--Penrose inverse $B^\dagger$ that other rotor counts use. Given a commanded collective thrust $T$ and body-torque command $\bm u$, the corresponding inverse-allocation and clipping model is
\begin{equation}
 \begin{aligned}
 \bm f_{\rm req}&=B^{-1}[T,\bm u\T]\T,\\
 \bm f_{\rm act}&=\clip_{[\bm0,\bm f_{\max}]}(\bm f_{\rm req}),\\
 [T_{\rm act},\bm\tau_{\rm act}\T]\T&=B\bm f_{\rm act}.
 \end{aligned}
 \label{eq:mixer}
\end{equation}
Here $\bm u=\bm u_0$ for PD; the CMG path in Section~\ref{sec:architecture}
modifies that requested torque before the same allocation and clipping.
Each $f_i$ is a commanded force in newtons, not a rotor-speed command.
Exact wrench realization requires
$B^{-1}[T,\bm u\T]\T\in[\bm0,\bm f_{\max}]$.
A componentwise torque box is not generally contained in this
collective-dependent feasible set. Clipped inverse allocation is used,
not an optimization-based allocator \cite{monteiro2016allocation}.

Collective thrust remains classical, with $z_d=1$ m:
\begin{equation}
 T=\clip_{[0,2.5Mg]}
 \left(\frac{M[g+6(z_d-z)-4\dot z]}{\max(R(Q)_{33},0.35)}\right).
 \label{eq:collective}
\end{equation}
Here $g$ is gravity; the height loop reads altitude and vertical velocity,
and the denominator compensates tilt. The guard $0.35$ avoids near-singular
compensation as $R(Q)_{33}\to0$, limiting it to about
$\arccos(0.35)\approx69.5^\circ$; the upper clip does not supersede
individual rotor limits. The CMG path has no direct collective
output, but can affect height through attitude, tilt compensation, and
allocation saturation.
The resulting shared plant interface isolates changes in requested
attitude torque while preserving the same physical simulation for both
controllers.

%%%%%%%%%%%%%%%%%%%%%%%%%%%%%%%%%%%%%%%%%%%%%%%%%%%%%%%%%%%%%%%%%
\section{CMG State Dynamics and Online Decoder Adaptation}
\label{sec:architecture}
\label{sec:analysis}
%%%%%%%%%%%%%%%%%%%%%%%%%%%%%%%%%%%%%%%%%%%%%%%%%%%%%%%%%%%%%%%%%
This section specifies the residual architecture in state-space form and the online decoder update, then states the boundedness properties that follow from the resulting saturations and projections.
\subsection{Encoding, States, and Numerical Update}

For $m$ torque channels and $2m$ units, normalize the clipped baseline and
form positive/negative current pairs:
\begin{equation}
 \begin{aligned}
 \bar{\bm u}&=\bm u_0\oslash\bm\ell,\\
 \bm I^{\pm}&=I_b\bm1+I_g\max(\pm\bar{\bm u},\bm0),\\
 \bm I&=[(\bm I^+)\T,(\bm I^-)\T]\T.
 \end{aligned}
 \label{eq:encoding}
\end{equation}
Each unit has a dimensionless membrane state $v_j$ and gate $r_j$:
\begin{align}
 \dot r_j&=[\sigma(\alpha_h(v_j-v_{\rm th}+\gamma r_j))-r_j]/\tau_r,\\
 \dot v_j&=-(v_j-v_{\rm rest})/\tau+I_j/\tau-c_j(v_j-v_{\rm reset}),\\
 c_j&=(g_0+\chi g_1|I_j|)r_j^2,\qquad \chi\in\{0,1\}.
 \label{eq:cmg}
\end{align}
% [v10] Coupling-topology figure cut for the 8-page target (fully specified
% already by \eqref{eq:cmg}); TikZ source kept under architectures/tikz/
% in case a journal/extended version restores it.
The logistic gate target depends on $v_j$, while $r_j$ shifts that target
and modulates membrane conductance: the membrane sets the gate target, and
the gate acts back twice, shifting its own threshold through $\gamma$ and
shunting the membrane through $c_j\propto r_j^2$ ($\chi=1$ makes that shunt
input dependent). Thus, both states participate even
though the decoder reads the gate-derived activity. The model has no
threshold-triggered reset or spike event; $v_{\rm reset}$ is a reversal
level in a continuous conductance term. The gains $g_0,g_1$ have the units
needed to make $c_j$ a rate.

The actual integration is sequential, not exponential Euler for both
states. From the old states compute
\begin{align}
 b_{j,k}&=\sigma\!\left(\clip_{[-60,60]}
       [\alpha_h(v_{j,k}-v_{\rm th}+\gamma r_{j,k})]\right),\\
 r_{j,k}^{+}&=\clip_{[0,1]}
       [r_{j,k}+(h/\tau_r)(b_{j,k}-r_{j,k})].
 \label{eq:gate_step}
\end{align}
Then evaluate $c_{j,k}=(g_0+\chi g_1|I_{j,k}|)(r_{j,k}^{+})^2$,
$\lambda_{j,k}=\tau^{-1}+c_{j,k}$, and
$d_{j,k}=(v_{\rm rest}+I_{j,k})/\tau+c_{j,k}v_{\rm reset}$. The
frozen-coefficient membrane step is
\begin{equation}
 v_{j,k}^{+}=e^{-\lambda_{j,k}h}v_{j,k}
       +(1-e^{-\lambda_{j,k}h})d_{j,k}/\lambda_{j,k}.
 \label{eq:mem_step}
\end{equation}
This exactly solves the scalar membrane equation with those frozen
coefficients; it does not establish unconditional stability of the coupled
nonlinear closed loop. Assign $(v_{j,k+1},r_{j,k+1})=(v_{j,k}^{+},r_{j,k}^{+})$.
The exponential-Euler step
$r_{j,k}^{+}=e^{-h/\tau_r}r_{j,k}+(1-e^{-h/\tau_r})b_{j,k}$, exact for a
frozen target, is convex for every $h>0$ without a clip
\cite{hochbruck2010exponential}; it changes only the weight on $r_{j,k}$
($0.905$ vs.\ $0.9$ at $h/\tau_r=0.1$), and all results use
\eqref{eq:gate_step}.

\subsection{Readout, PES, and Additive Residual}
The within-step activity, initial decoder, and clipped readout are
\begin{equation}
 \begin{aligned}
 \bm a_k&=\clip_{[0,1]}(G\bm r_k^+),\qquad W_0=[I_m,-I_m],\\
 \bm s_k&=\clip_{[-1,1]}(W_k\bm a_k).
 \end{aligned}
 \label{eq:readout}
\end{equation}
Pair subtraction retains directional information from a nonnegative
activity basis; state memory does not imply instantaneous sign agreement
after a command reversal. With $G=10^5$ the activity saturates once
$r_j\ge10^{-5}$, so $G$ is a primary design hyperparameter for how much of
$r_j$'s dynamic range reaches the decoder. The teaching target and local update are
\begin{align}
 \bm y_k^*&=\clip_{[-1,1]}(-k_e\bm e_{q,k}-k_\omega^*\bm\omega_k),\\
 \bm\varepsilon_k&=\bm y_k^*-\bm s_k,\\
 W_{k+1}&=\Pi_{\mathcal W}
       [W_k+\eta\bm\varepsilon_k\bm a_k\T-\lambda_W W_k].
 \label{eq:pes}
\end{align}
Here $\mathcal W=[-W_{\max},W_{\max}]^{m\times2m}$. The coefficients $\eta$ and $\lambda_W$ are per-step (so changing $h$ changes the effective learning dynamics). PES is a local error-driven update toward the prescribed target $\bm y_k^*$; it does not differentiate through the plant or claim to minimize a closed-loop cost. Since $\bm s_k$ is normalized by the authority $\beta\bm\ell$, anisotropic scaling enters only through $\bm\ell$ and one pair $(k_e,k_\omega^*)$ serves all axes.

Using $W_k$ before the update, the requested command is
\begin{equation}
 \bm u_k=\sat_{\bm\ell}(\bm u_{0,k}+\beta\bm\ell\odot\bm s_k),\quad
 \bm\delta_k=\bm u_k-\bm u_{0,k},\quad 0\le\beta\le1.
 \label{eq:residual}
\end{equation}
The canonical UAV case has $m=3$, six CMG units, and eighteen decoder
weights. Intrinsic conductance modulation ($\chi$), synaptic/decoder
adaptation ($W_k$), and output authority ($\beta\bm\ell$) are distinct.

\subsection{Structural Properties and Their Scope}
We first state the boundedness properties that follow directly from the architecture: internal-state bounds, a uniform command-deviation bound set by residual authority, and projection-based bounds on decoder weights.
\begin{proposition}[Bounded internal states]
\label{prop:states}
With the nonnegative conductance parameters in Table~\ref{tab:parameters}
and $r_{j,0}\in[0,1]$, the gate remains in $[0,1]$ for every $h>0$, and the
gate clip in \eqref{eq:gate_step} is inactive when $0<h\le\tau_r$.
The membrane iterate remains in the convex hull of its initial value,
$v_{\rm reset}$, and the interval
$[v_{\rm rest}+I_b,v_{\rm rest}+I_b+I_g]$.
\end{proposition}
\begin{proof}
For $0<h\le\tau_r$ the unclipped gate step is a convex combination of
$r_{j,k}$ and $b_{j,k}\in[0,1]$, so the clip is inactive; for $h>\tau_r$ the
clip enforces the bound. Since $I_j\in[I_b,I_b+I_g]$ and $c_j\ge0$,
$d_j/\lambda_j$ is a convex combination of $v_{\rm rest}+I_j$ and
$v_{\rm reset}$. Equation~\eqref{eq:mem_step} is a convex combination of
$v_{j,k}$ and $d_j/\lambda_j$. Induction proves both statements.
\end{proof}

\begin{proposition}[Bounded authority and baseline recovery]
\label{prop:authority}
For $\bm u_{0,k}\in\mathcal U$ and any finite $W_k$,
$\bm u_k\in\mathcal U$ and $|\delta_{k,i}|\le\beta\ell_i$.
If $\beta=0$, the requested baseline command is recovered at every step.
If $W_k=0$ at a particular step, it is recovered at that step, but learning
need not keep $W_{k+1}=0$.
\end{proposition}
\begin{proof}
The outer clip enforces membership in $\mathcal U$. Its componentwise
nonexpansiveness and $\sat_{\ell_i}(u_{0,k,i})=u_{0,k,i}$ give
$|\delta_{k,i}|\le|\beta\ell_i s_{k,i}|\le\beta\ell_i$.
The recovery statements follow by substitution.
\end{proof}
With $\beta=0$ the same allocator and state give the same actuator output.
The proposition establishes neither state safety, rotor headroom, stability,
nor a bound on tracking improvement; a bounded adversarial correction can
still degrade or destabilize a nonlinear plant.

\begin{proposition}[Decoder bound and restricted contraction]
\label{prop:pes}
If $W_0\in\mathcal W$, then \eqref{eq:pes} keeps $W_k\in\mathcal W$.
If, additionally, the target and a nonzero activity $\bm a$ are fixed,
leak is zero, and readout clipping and projection remain inactive, then
\begin{equation}
 \bm\varepsilon_{k+1}
       =(1-\eta\|\bm a\|_2^2)\bm\varepsilon_k.
 \label{eq:contraction}
\end{equation}
For a nonzero initial error, its norm contracts to zero when
$0<\eta\|\bm a\|_2^2<2$.
\end{proposition}
\begin{proof}
Projection gives the first result. Under the additional conditions,
$(W_{k+1}-W_k)\bm a=\eta\bm\varepsilon_k\|\bm a\|_2^2$;
subtracting this from the fixed target error gives
\eqref{eq:contraction}. The scalar multiplier has modulus less than one
under the stated condition.
\end{proof}
Beyond those conditions, the readout clip is a nonexpansive projection onto
a convex set containing $\bm y_k^*$ \cite{bauschke2017convex}, so
$\|\clip_{[-1,1]}(\bm z)-\bm y_k^*\|\le\|\bm z-\bm y_k^*\|$: saturation never
enlarges the PES error.
This is the constant-feature PES result \cite{voelker2015pes}, not a
closed-loop convergence theorem. Since $\|\bm a\|^2\le2m$,
$0<\eta<1/m$ is sufficient for that restricted model with nonzero
activity; changing targets, loss of excitation, clipping, and leakage
invalidate a direct application to the experiment. No unique identification
of all decoder weights follows from one fixed regressor.

% [v7] The statement a control audience looks for first. Scoped: it is a
% conditional bound, not a stability theorem, and says so.
%\begin{remark}[Scope of a small-gain argument]
%\label{rem:smallgain}
%Propositions~\ref{prop:authority} and \ref{prop:pes} place the realized
%command deviation in a known box, $|\delta_{k,i}|\le\beta\ell_i$, uniformly
%in $k$ and independently of what the decoder has learned. If the baseline
%loop is input-to-state stable with respect to an additive torque
%perturbation, with gain $\gamma(\cdot)$ on that input, then the augmented
%loop inherits an ISS estimate with asymptotic bound
%$\gamma(\beta\|\bm\ell\|)$, and $\beta$ tunes it continuously to zero. Two
%conditions are not discharged here. The perturbation entering the plant is
%post-allocation, so the relevant bound is on $B^{-1}$ applied to
%$\bm\delta_k$ under rotor saturation, not on $\bm\delta_k$ itself; and the
%required ISS gain for the saturated quaternion-PD baseline is not
%established in this paper. The architecture therefore supplies the bounded
%perturbation such an argument needs, and identifies the two missing
%ingredients, without asserting the conclusion.
%\end{remark}

\subsection{Allocation-Aware Local Stability}
\label{sec:local}
Let $E=[\bm0_{3\times1},I_3]\T$, $H=E\T$, and let
$\mathcal A_T(\bm u)=HB\clip_{[0,\bm f_{\max}]}(B^{-1}[T,\bm u\T]\T)$
be the realized torque at collective $T$. For two commands at the
\emph{same state and collective}, scalar clipping has secant slopes
$d_i\in[0,1]$. Consequently,
\begin{align}
 \mathcal A_T(\bm u)-\mathcal A_T(\bm u_0)
 &=HBDB^{-1}E\bm\delta,\\
 \|\mathcal A_T(\bm u)-\mathcal A_T(\bm u_0)\|_2
 &\le L_A\beta\|\bm\ell\|_2,
 \label{eq:allocbound}
\end{align}
where $D=\operatorname{diag}(d_i)$ and
$L_A=\max_{D\in\mathcal D}\|HBDB^{-1}E\|_2$.
$\mathcal D$ is the set of sixteen diagonal matrices with entries
in $\{0,1\}$. The vertex maximum bounds all intermediate $D$ by convexity.
Numerically $L_A=\certXLA$ (X2) and $\certCLA$ (CF2), inflated because
small $|\kappa_i|$ turn clipped yaw requests into roll/pitch torque. If all
of $\mathcal U$ is rotor-feasible at the operating collective (CF2 at hover,
not X2), $D=I$ and $1$ replaces $L_A$.

For a constant reference and a smooth shortest-path branch, let
$\bm x=[\bm e_q\T,\bm\omega\T]\T$ and
$K_q=\operatorname{diag}(\bm k_q)$, $K_\omega=\operatorname{diag}(\bm k_\omega)$.
Write the continuous-time attitude subsystem as
\begin{equation}
 \begin{gathered}
 \dot{\bm x}=A\bm x+\bm g(\bm x)+F\bm d,\\
 A=\begin{bmatrix}0&\frac12I_3\\-\bm J^{-1}K_q&-\bm J^{-1}K_\omega\end{bmatrix},
 \qquad F=\begin{bmatrix}0\\\bm J^{-1}\end{bmatrix},
 \end{gathered}
 \label{eq:localmodel}
\end{equation}
where $\bm d=\bm\tau_{\rm act}+\bm\tau_{\rm p}+\bm\tau_{\rm d}
 +K_q\bm e_q+K_\omega\bm\omega$ is the total deviation from ideal PD torque.
It includes allocation and command-saturation errors, not just the neural
request. The nonlinear remainder is
\begin{equation}
 \bm g(\bm x)=\begin{bmatrix}
 \tfrac12[(\sqrt{1-\|\bm e_q\|^2}-1)I_3+[\bm e_q]_\times]\bm\omega\\
 -\bm J^{-1}[\bm\omega]_\times\bm J\bm\omega
 \end{bmatrix}=O(\|\bm x\|^2),
 \label{eq:localremainder}
\end{equation}
with $[\bm z]_\times\bm y=\bm z\times\bm y$; norms use fixed coordinates,
rates in rad/s. Bounding each block separately (orthogonality and
$\|\bm e_q\|\|\bm\omega\|\le\frac12\|\bm x\|^2$ for the first; the identity
$\bm\omega\times\bm J\bm\omega=\bm\omega\times(\bm J-\mu I_3)\bm\omega$ with
$\mu=\frac12(\lambda_{\max}+\lambda_{\min})(\bm J)$ for the second, $\kappa$
the condition number; full steps in the companion report) gives
\begin{equation}
 \|\bm g(\bm x)\|\le L_r\|\bm x\|^2,\qquad
 L_r=\tfrac{\sqrt2}{4}+\tfrac12(\kappa(\bm J)-1),
 \label{eq:Lr}
\end{equation}
i.e., $L_r=\certXLr$ (X2) and $\certCLr$ (CF2), above sampled suprema
$\certXLrEmp$ and $\certCLrEmp$.

\begin{theorem}[Conditional local practical stability]
\label{thm:local}
Let $\bm J,K_q,K_\omega\succ0$, let $q_d$ be constant, and use the smooth
branch in \eqref{eq:localmodel}. Choose $Q_L\succ0$ and the solution
$P\succ0$ of $A\T P+PA=-Q_L$. On $\|\bm x\|\le r<1$, suppose
$\|\bm g(\bm x)\|\le L_r\|\bm x\|^2$ (e.g., \eqref{eq:Lr}) and choose $r$ so that
$c=\lambda_{\min}(Q_L)-2\|P\|L_r r>0$.
Set $p_- =\lambda_{\min}(P)$, $p_+=\lambda_{\max}(P)$,
$b=\|PF\|$, and $V=\bm x\T P\bm x$. If
\begin{equation}
 V(0)\le p_-r^2,\qquad
 \|\bm d(t)\|\le\bar d<\frac{cr}{2b}\sqrt{\frac{p_-}{p_+}},
 \label{eq:localsmallness}
\end{equation}
then $V\le p_-r^2$ is forward invariant and
\begin{equation}
 V(t)\le e^{-\alpha t}V(0)+\frac{4p_+b^2}{c^2}\bar d^2,
 \qquad \alpha=\frac{c}{2p_+}.
 \label{eq:localiss}
\end{equation}
Thus attitude regulation is locally input-to-state stable; zero
perturbation gives local exponential convergence of the plant state.
\end{theorem}
\begin{proof}
Along $\dot{\bm x}=A\bm x$, the mechanical energy
$V_L=\bm e_q\T K_q\bm e_q+\frac12\bm\omega\T\bm J\bm\omega$ satisfies
$\dot V_L=\bm e_q\T K_q\bm\omega-\bm\omega\T K_q\bm e_q
-\bm\omega\T K_\omega\bm\omega=-\bm\omega\T K_\omega\bm\omega\le0$ because
$K_q$ is symmetric; off-diagonal products of inertia in $\bm J$ do not
enter. On the set $\dot V_L\equiv0$, $\bm\omega\equiv\bm0$ forces
$\dot{\bm\omega}\equiv\bm0$ and hence $\bm e_q=\bm0$, so by LaSalle's
invariance principle \cite{khalil2002nonlinear} $A$ is Hurwitz and the
unique $P\succ0$ exists.
Equations~\eqref{eq:localmodel}--\eqref{eq:localremainder} give
\[
 \dot V\le-c\|\bm x\|^2+2b\|\bm x\|\bar d
 \le-\frac c2\|\bm x\|^2+\frac{2b^2}{c}\bar d^2.
\]
On $V=p_-r^2$, $\|\bm x\|\ge r\sqrt{p_-/p_+}$, so
\eqref{eq:localsmallness} makes the first bound strictly negative.
The ellipsoid stays inside the radius-$r$ ball; integration of
$\dot V\le-\alpha V+2b^2\bar d^2/c$ gives \eqref{eq:localiss}.
\end{proof}

If the same-state baseline deviation from ideal PD is bounded by $\bar d_0$,
then \eqref{eq:allocbound} permits
$\bar d=\bar d_0+L_A\beta\|\bm\ell\|$ in the theorem, \emph{provided}
\eqref{eq:localsmallness} holds, independent of the particular PES
trajectory. Maximizing the admissible bound in
\eqref{eq:localsmallness} over $r$ gives
$r^*=\min\{\lambda_{\min}(Q_L)/(4\|P\|L_r),1\}$ and, when $r^*<1$,
\begin{equation}
 \bar d_{\max}=\frac{\lambda_{\min}^2(Q_L)}{16\,b\,\|P\|\,L_r}
 \sqrt{\frac{p_-}{p_+}},\qquad
 \beta<\beta_{\max}=\frac{\bar d_{\max}-\bar d_0}{L_A\|\bm\ell\|_2},
 \label{eq:betamax}
\end{equation}
a sufficient authority bound when $\bar d_0<\bar d_{\max}$. With $Q_L=I$
(no $\operatorname{diag}(q_eI_3,I_3)$ did better), $A$ has spectral abscissa
$\certXEigA$ (X2) and $\certCEigA$ (CF2), and the PD command is unsaturated
and rotor-feasible in the certified ball, so $\bar d_0=0$ at hover. Then
$(r^*,\bar d_{\max},\beta_{\max})=(\certXrstar,\ \certXdbarmax\,\mathrm{N\,m},\
\certXbetamax)$ for X2 and $(\certCrstar,\ \certCdbarmax\,\mathrm{N\,m},\
\certCbetamax)$ for CF2 ($\certCbetamaxNoClip$ with exact allocation). The
experimental $\beta=0.5$ lies far outside this conservative certificate:
stability of the reported runs is observed, not certified. This bound
holds under rotor clipping but not between two different closed-loop
trajectories, and is not a PES convergence proof or a height-error bound;
moving references, the $180^\circ$ sign boundary, sampled-data/solver
error, and full translational stability remain outside its scope.

\begin{remark}[Sign selection on certified trajectories]
Under Theorem~\ref{thm:local}, $\|\bm e_q\|\le\|\bm x\|\le r<1$ gives
$|q_{e,0}|\ge\sqrt{1-r^2}>0$: sign selection in \eqref{eq:theta} never
switches, the smooth branch is self-consistent, and $q_{e,0}$ perturbations
below $\sqrt{1-r^2}$ cannot flip it; the global obstruction
\cite{bhat2000topological} remains.
\end{remark}

\paragraph{Nengo realization.}
A custom Nengo~4.1.0 \texttt{NeuronType} implements the CMG states directly, with matched PES hyperparameters ($\eta_N=\eta(2m)/h$) and equivalent leak/projection handling; a one-step teaching-graph latency makes the weight sequences differ slightly from the direct backend. Over ten matched cases the maximum relative mean-error difference is $2.013\times10^{-6}$: a tested numerical agreement, not formal backend equivalence.


\section{Experimental Protocol and Cross-Mission Results}
\label{sec:results}
\subsection{Models, Parameters, and Measurements}

MuJoCo \cite{todorov2012mujoco} simulates two Menagerie-derived multirotor models \cite{menagerieModels}, CF2 and X2, as free-flying rigid bodies with site-applied wrenches; the attitude loop commands three body torques and a separate height loop commands collective thrust.
The upstream CF2 model exposes abstract moment channels \cite{menagerieCF2}; we replace them with four rotor-force channels and the shared clipped inverse allocator \eqref{eq:mixer} to expose thrust/torque coupling and per-rotor limits. Rotor ordering is $(-x,-y),(-x,+y),(+x,+y),(+x,-y)$ with reaction signs $(-,+,-,+)$. The CF2 $1.5\times$ scale (Table~\ref{tab:models}) raises a model limit, not measured hardware capability. Both models are simplified (no motor/RPM/battery states and no validated rotor aerodynamics); feedback uses \texttt{qpos}/\texttt{qvel}, with noise injected only in Sec.~\ref{sec:noise}.

\begin{table}[t]
\centering
\caption{Model/interface parameters with manufacturer context (top) and
canonical 3-channel CMG--PES parameters (bottom). Mass agrees with published CF2/X2 specifications \cite{bitcrazeCF20Specs,skydioX2ESpecs}; force limits, inertias, and reaction coefficients are simulation parameters, not identified hardware ratings. Neural states/currents are normalized model variables; learning coefficients are per step.}
\label{tab:models}\label{tab:parameters}
\setlength{\tabcolsep}{3pt}
\footnotesize
\begin{tabular}{@{}p{0.33\linewidth}p{0.30\linewidth}p{0.29\linewidth}@{}}
\toprule
Parameter & CF2 rotor model & X2 model\\
\midrule
Mass $M$ (kg) & $0.027$ & $1.325$\\
Rotor $|p_x|,|p_y|$ (m) & $0.0325,0.0325$ & $0.14,0.18$\\
Reaction $|\kappa_i|$ (m) & $0.005$ & $0.0201$\\
Rotor/total force limit (N) & $.0875\to.13125$ / $.350\to.525$ & $13$ / $52$\\
Main-test thrust/weight & $1.982$ & $4.001$\\
$\bm j_{\rm tune}$ (kg\,m$^2$) & $10^{-5}(2.3951,\allowbreak2.3951,\allowbreak3.2347)$ & $10^{-2}(6.0711,\allowbreak3.6468,\allowbreak2.5412)$\\
Integrator, timestep & RK4, $0.0005$ s & Euler, $0.0005$ s\\
Mfr. thrust data & $.568$ N\textsuperscript{a} & Not published\textsuperscript{b}\\
\bottomrule
\end{tabular}
\par\smallskip
\begin{minipage}{0.97\linewidth}
\scriptsize
\textsuperscript{a}Historical CF2 measurement at 93.75\% PWM, not a continuous rating \cite{bitcrazeThrustMeasurements}; the modeled $.525$ N is below it.
\textsuperscript{b}Not reported for X2; neither model includes motor lag (CF2's historical step response is $\sim$180\,ms \cite{bitcrazeMotorResponse}).
\end{minipage}
\par\medskip
\footnotesize\setlength{\tabcolsep}{3pt}
\begin{tabular}{@{}ll@{\hspace{8pt}}ll@{}}
\toprule
CMG--PES & Value & & Value\\\midrule
$h$, duration & $0.0005$ s, $45$ s & $\alpha_h,\gamma,g_0,g_1$ & $50,2,120,30$\\
$\tau,\tau_r$ & $0.05$ s, $0.005$ s & $I_b,I_g,G$ & $0.75,1.25,10^5$\\
$v_{\rm rest},v_{\rm reset},v_{\rm th}$ & $0,0,1$ & $\beta,\eta$ & $0.5,8\times10^{-4}$\\
$v_{j,0},r_{j,0},\chi$ & $0,0,1$ & $\lambda_W,W_{\max},k_e,k_\omega^*$ & $10^{-6},6,8,0$\\
\bottomrule
\end{tabular}
\end{table}

\begin{comment}
\begin{table}[t]
\centering
\caption{Simulation-model characteristics. CF2 authority is scaled relative to the upstream model limit.}
\label{tab:models}
\setlength{\tabcolsep}{3pt}
\footnotesize
\begin{tabular}{@{}p{0.40\linewidth}p{0.27\linewidth}p{0.25\linewidth}@{}}
\toprule
Property & CF2 rotor model & X2 model\\
\midrule
Mass $M$ (kg) & $0.027$ & $1.325$\\
Rotor $|p_x|,|p_y|$ (m) & $0.0325,0.0325$ & $0.14,0.18$\\
Reaction magnitude $|\kappa_i|$ (m) & $0.005$ & $0.0201$\\
Unscaled rotor limit (N) & $0.0875$ & $13$\\
Main-test authority scale & $1.5$ & $1$\\
Main-test rotor limit (N) & $0.13125$ & $13$\\
Main-test total limit (N) & $0.525$ & $52$\\
Maximum thrust/weight & $1.321\to1.982$ & $4.001$\\
Rigid-body integrator & RK4 & Euler\\
Experiment timestep (s) & $0.0005$ & $0.0005$\\
Motor activation/RPM states & None & None\\
\bottomrule
\end{tabular}
\end{table}

\begin{table}[t]
\centering
\caption{Simulation parameters vs. manufacturer information (where available).}
\label{tab:manufacturer}
\setlength{\tabcolsep}{3pt}
\footnotesize
\begin{tabular}{@{}>{\raggedright\arraybackslash}p{0.25\linewidth}>{\raggedright\arraybackslash}p{0.30\linewidth}>{\raggedright\arraybackslash}p{0.37\linewidth}@{}}
\toprule
Quantity & Simulation & Manufacturer information\\
\midrule
CF2 mass & $27$ g & $27$ g \cite{bitcrazeCF20Specs}\\
X2 mass & $1325$ g & $1325$ g with battery \cite{skydioX2ESpecs}\\
CF2 total thrust & $0.350\to0.525$ N & $0.568$ N measured at $93.75\%$ PWM\textsuperscript{a} \cite{bitcrazeThrustMeasurements}\\
X2 rotor limit & $13$ N per rotor & Not reported \cite{skydioX2ESpecs}\\
CF2 motor response & No motor dynamics & $\sim180$ ms: $0\to25{,}000$ RPM\textsuperscript{b} \cite{bitcrazeMotorResponse}\\
Inertia tensors & CF2 inherited; X2 geometry-derived & Not reported in product specifications \cite{bitcrazeCF20Specs,skydioX2ESpecs}\\
\bottomrule
\end{tabular}
\par\smallskip
\begin{minipage}{0.97\linewidth}
\scriptsize
\textsuperscript{a}Historical total-thrust measurement with a power-supply test setup, not a guaranteed continuous battery-powered rating.
\textsuperscript{b}One motor's reported full-command response with a charged battery; not an identified first-order time constant.
\end{minipage}
\end{table}
\end{comment}

CF2's body axes are principal axes. X2's geometry-derived body tensor,
$\bm J=10^{-2}[3.66517,0,-.21000;\allowbreak\,0,2.54117,0;\allowbreak\,-.21000,0,6.05280]$\,kg\,m$^2$
(rows separated by semicolons), differs from both
$\operatorname{diag}(\bm j_{\rm tune})$ and a manufacturer-identified
tensor. Torque limits use $\bm j_{\rm tune}$ as defined in Section~\ref{sec:problem}.

Each PD/CMG pair shares initial state, reference, duration, allocator, model, and timestep. The primary campaign is nominal and deterministic (no injected disturbances or sensor noise), one run per cell, learning active throughout. We report the mean error $\bar\theta_q$ and the reduction $100(\bar\theta_{q,\rm PD}-\bar\theta_{q,\rm CMG})/\bar\theta_{q,\rm PD}$ over all 90{,}000 samples, even when plotted traces are decimated.

The five reference triples are generated as follows (radians):
Roll is $(0.25\sin0.7t,0,0)$; Full3 is
$0.5\sin(t+(0,\pi/2,\pi)\T)$, componentwise; Hard45 is
$(0.55\sin0.75t+0.20\sin(2.10t+0.40),\allowbreak\;
0.45\sin(0.95t+\pi/3)+0.18\sin1.75t,\allowbreak\;
0.65\sin(0.55t+\pi)+0.25\sin(1.35t+0.20))$.
Hardstep holds zero for one second, then cycles every two seconds through
$(.75,-.45,.55)$, $(-.65,.55,-.70)$, $(.50,.50,.90)$, and $(-.80,-.35,.20)$.
Yaw360 is $(0,0,2\pi(3\xi^2-2\xi^3))$ with $\xi=t/45$.
Hover is omitted: its near-zero error admits no meaningful percentage.

\begin{figure}[t]
\centering
\includegraphics[width=\linewidth]{figures/02_neural_activity/attitude_detail/x2_full3_attitude.png}
\caption{X2 Full3 mission: PD baseline versus CMG--PES augmentation over the
45-s run. Full3 commands all three axes simultaneously and is the largest
reduction that is also insensitive to measurement noise, losing under one
point between the clean and realistic columns of Table~\ref{tab:noise}.}
\label{fig:tracking}
\end{figure}

% [v10] Cut for the 8-page target: this panel duplicated the mechanistic
% narrative already given in full in Sec.~IV-C (Ablation); the membrane/
% gate/activity traces are in the companion derivations report.

\subsection{Primary Five-Mission Matrix}
CMG-PES improves mean quaternion-angle tracking error in all ten nominal missions while respecting the same actuator interface as the PD baseline.
%
\begin{table}[h]
\caption{Canonical Nengo CMG--PES versus PD: mean quaternion-angle error
[rad] and reduction vs PD [\%], $45$ s, $n=1$ per clean cell. PES: CMG--PES;
Fixed: same units with $W_k=W_0$; Mod./Real.: CMG--PES under the moderate
($0.010$\,rad, $0.025$\,rad\,s$^{-1}$) and realistic ($0.020$\,rad,
$0.050$\,rad\,s$^{-1}$) attitude/rate noise of Sec.~\ref{sec:noise}, two-seed
means with seed SD $\le0.03$ pp at the realistic level. CF2 denotes the
$1.5\times$ rotor-authority design, not stock hardware.}
\label{tab:mujoco}\label{tab:noise}\centering\scriptsize
\setlength{\tabcolsep}{2.4pt}
\begin{tabular}{@{}llrrrrrr@{}}
\toprule
&&\multicolumn{2}{c}{Error [rad]}&\multicolumn{4}{c}{Reduction vs PD [\%]}\\
\cmidrule(lr){3-4}\cmidrule(l){5-8}
Model&Mission&PD&CMG--PES&PES&Fixed&Mod.&Real.\\\midrule
X2&Roll&.050565&.041785&17.36&1.93&32.3&32.2\\
&Full3&.303275&.131748&56.56&9.74&56.4&55.7\\
&Hard45&.294082&.138007&53.07&10.27&52.9&52.3\\
&Hardstep&.454528&.359355&20.94&0.46&20.8&20.6\\
&Yaw360&.075863&.029902&60.58&0.12&56.8&47.8\\\midrule
CF2&Roll&.051337&.042439&17.33&3.31&32.4&32.5\\
&Full3&.288383&.126476&56.14&4.53&55.9&55.2\\
&Hard45&.285449&.132319&53.65&6.80&53.4&52.7\\
&Hardstep&.426182&.347886&18.37&0.42&18.3&18.1\\
&Yaw360&.075812&.029870&60.60&0.06&56.8&47.8\\\bottomrule
\end{tabular}
\end{table}
%
The largest reductions occur on Yaw360 and Full3 and the smallest on Roll.
Hardstep has the largest absolute mean error and the only nonzero rotor-bound
occupancy $S$, the logged percentage of samples with at least one rotor
command at a bound (PD/CMG): $.016/.239$ on X2 and $.137/.264$ on CF2. The
statistic $S$ measures post-clipping commands near a limit, not
requested-wrench infeasibility magnitude. The results
establish deterministic nominal performance relative to this fixed PD law,
not optimality over classical controllers.

Fig.~\ref{fig:tracking} shows the X2 Full3 tracking behind these numbers;
the paired membrane/gate/activity traces are discussed in Sec.~\ref{sec:ablation}
and given in full in the companion report. Improvement refers to the mean
metric; it need not hold at every instant.


%The primary five-mission matrix above shows that all ten cases reduce mean quaternion-angle error under CMG--PES. Gains are largest on Yaw360 and Full3 and smallest on Roll; the abrupt Hardstep case is intermediate. The reported bound-occupancy statistic detects post-clipping commands within the runner's tolerance of a limit; it is not a measure of requested-wrench infeasibility magnitude. Hardstep has small but nonzero occupancy on both models, so a blanket zero-saturation claim would be incorrect. These cases establish repeatable nominal performance; robustness to model mismatch and comparison against other baseline families are separate questions, addressed in Sec.~\ref{sec:conclusion}.

%The two featured cases bracket the matrix. Figures~\ref{fig:tracking} and \ref{fig:full3_neurons} show X2 Full3, the largest reduction that is also noise-insensitive, with the neural activity that produced it; Figs.~\ref{fig:states} and \ref{fig:cf2_hardstep_neurons} show CF2 Hardstep, the smallest reduction and the only case whose rotor commands reach a bound. Yaw360 carries the largest clean reduction but is also the most noise-sensitive entry in Table~\ref{tab:noise}, so it is reported rather than featured.

% [v9] Noise sensitivity: formerly a separate table; its entries are the
% Mod./Real. columns of tab:mujoco (regenerate from
% results/final_sweep_2026-09-05/compiled/table_noise.tex). Per-mission
% realistic seed SDs: Roll .01, Full3 .00, Hard45 .01, Hardstep .02, Yaw360 .03 pp.

\begin{comment}
\begin{figure}[t]
\centering
\includegraphics[width=\linewidth]{figures/02_neural_activity/attitude_detail/cf2_rotor_hardstep_attitude.png}
\caption{CF2 Hardstep mission: PD baseline versus CMG--PES augmentation on
the abrupt step reference. This is the hardest case in the matrix, the
smallest reduction ($18.4\%$), and the only one whose rotor commands reach a
bound ($S=.137/.264$). The residual stays inside its
authority envelope and the augmented loop remains better than the baseline
throughout.}
\label{fig:states}
\end{figure}
\end{comment}

%%%%%%%%%%%%%%%%%%%%%%%%%%%%%%%%%%%%%%%%%%%%%%%%
\subsection{Sensitivity to Measurement Noise}
\label{sec:noise}

Each of the ten cases was rerun at the two noise levels of
Table~\ref{tab:noise} with two paired seeds: at every step, independent
zero-mean Gaussian samples form a small Euler-angle quaternion
left-multiplied into $q$ and corrupt each angular-rate component. These
levels are not calibrated sensor specifications, and the runner evaluates
$\theta_q$ from the \emph{corrupted} quaternion used by the controller, not
a ground-truth attitude, so the Mod./Real. columns measure noisy-feedback
error. All noisy reductions remain positive, spanning $18.08$--$55.75\%$ at
the higher level; Yaw360 loses about 12.8 pp relative to clean feedback,
whereas Roll's reduction increases. Seed SDs are $.001$--$.101$ pp
(moderate) and $.002$--$.031$ pp (realistic): limited repeatability, not a
population confidence interval.

\begin{comment}
\begin{figure}[t]
\centering
\includegraphics[width=\linewidth]{figures/02_neural_activity/cf2-rotor_hardstep_neurons.png}
\caption{CF2 Hardstep: internal CMG signals on the same trajectory as
Fig.~\ref{fig:states}. The plot separates membrane state, raw gate $r_j$, and
the clipped activity $a_j=\clip_{[0,1]}(Gr_j)$ presented to the decoder,
illustrating the readout saturation of Sec.~IV-D under the one mission that
also drives the rotors to a bound.}
\label{fig:cf2_hardstep_neurons}
\end{figure}
\end{comment}

%%%%%%%%%%%%%%%%%%%%%%%%%%%%%%%%%%%%%%%
\subsection{Ablation: What Drives the Improvement?}\label{sec:ablation}

We separate decoder adaptation (PES), input-dependent conductance
modulation ($\chi$), and readout gain ($G$). Toggling $\chi$, which
retains the base shunt $g_0r_j^2$, is numerically indistinguishable with
the decoder mode held fixed. \emph{Fixed versus adaptive decoder.} Holding $W_k=W_0$ keeps the neuron
dynamics but reduces the readout to $\bm a^+-\bm a^-$, a bounded, stateful
image of the PD request. It still helps (X2 $0.12$--$10.27\%$, CF2
$0.06$--$6.80\%$; Table~\ref{tab:mujoco}), so fixing the decoder does not
remove \emph{all} improvement, yet it recovers at most a fifth of the PES
gain in any cell, for the mechanistic reasons below.
\emph{Membrane and gate.} At $I_j=I_b$ the membrane rests at
$v_j\approx0.75<v_{\rm th}$, so
$r_j=\sigma(\alpha_h(I_b-v_{\rm th}))=\sigma(-12.5)\approx3.7\times10^{-6}$:
the gate is deep in its logistic tail, not closed. Once $r_j\ge1/2$,
$\gamma=2$ and $v_j\ge0$ (Proposition~\ref{prop:states}) give
$b_{j,k}\ge1/2$ in \eqref{eq:gate_step}, so an opened gate cannot close
(none did in the twenty runs), and the shunt holds $v_j$ near $0.09$.
\emph{Activity.} With $G=10^5$ the resting gate already gives
$a_j\approx0.37$, graded by the input; a latched gate pins $a_j=1$. With
the fixed decoder only the $\pm y$ pair latches on Full3 (both models), so
$a^+_y-a^-_y=0$ and the pitch residual is exactly zero after 1\,s while
the sub-threshold units carry roll and yaw. Under PES on X2,
$\bm a_k\equiv\bm1$ from $2.02$\,s (all gates latched by $2.4$\,s): with
constant $\bar{\bm a}$ and inactive projection, $\bm z_k=W_k\bar{\bm a}$
obeys the leaky extension of Proposition~\ref{prop:pes}'s zero-leak
constant-feature contraction \eqref{eq:contraction} (full recursion in the
companion report), i.e., a first-order filter of the teaching target with
pole $\pesPole$ (about $\pesTau$\,s) while the readout clip is inactive.
With $\|\bar{\bm a}\|^2=2m$ its DC gain is
$6\eta/(6\eta+\lambda_W)\approx0.9998$:
the latched residual is \emph{lagged proportional} action, not integral
action (which needs a unit pole). For $|k_ee_{q,i}|<1$ it adds
$\beta k_e\ell_i/k_{q,i}=\beta k_e\nu_i/(2\omega_{n,i}^2)$, i.e.,
$\stiffRatioRP\times$ (roll, pitch) and $\stiffRatioY\times$ (yaw) the PD
stiffness on both vehicles, behind a $\pesCorner$\,rad/s corner. The smooth
references are sums of constant-frequency components
\cite{sandoval2018instantaneous}; their largest instantaneous frequency,
$2.1$\,rad/s (Hard45), lies below the corner (gain $\lagHardGain$, lag
$\lagHardPhase^\circ$), while noise in $\bm e_q$ above it is attenuated.
\emph{Residual and decoder.} The residual RMS rises from $0.026$ to
$0.123$\,N\,m and the decoder moves $\|W_T-W_0\|_F=0.37$; on Roll, Full3,
Hard45, and Hardstep PES applies $4$--$12\times$ the fixed residual RMS on
both models. Where no gate opens (Roll, Yaw360; $r_j\le1.6\times10^{-3}$)
PES lifts the driven membranes toward threshold ($v_j\le0.87$) and uses
the graded tail. Decoder adaptation is therefore the principal lever: the
CMG dynamics supply a bounded basis, latched or graded, on which PES
shapes the correction. The leak $\lambda_W$ is the $\sigma$-modification of
robust adaptive control \cite{ioannou1984sigma,ioannou1996robust}: at a fixed
point of \eqref{eq:pes} with constant $\bar{\bm a}$ and inactive clip and
projection, $\eta\bm\varepsilon\|\bar{\bm a}\|^2=\lambda_W\bm s$, an exact
PES mismatch of $\pesLeakLatched\|\bm s\|$ latched and $\pesLeakTail\|\bm s\|$
at the resting tail.

\subsection{Neuron-Family Comparison}
The Nengo benchmark covers two models, five missions, and eleven controllers: PD and fixed/PES variants of CMG, LIF, ALIF, Izhikevich \cite{izhikevich2003}, and LIFRate, all sharing the bounded residual, saturation, and allocation interface; only the neuron dynamics and native-state-to-activity map differ.
Fig.~\ref{fig:mech_dashboard} shows two patterns. First, PES is the
dominant lever in every family: fixed decoders rarely match their PES
counterparts, even with more expressive dynamics, extending
Sec.~\ref{sec:ablation} to every basis. Second, lower error is not free:
some Izhikevich--PES cells reach the smallest error, but at an average
normalized residual RMS near 0.37 versus 0.12 for CMG--PES and with larger
decoder movement, so equal output caps do not imply equal realized effort.
Proposition~\ref{prop:pes} explains the flat rows: its factor
$1-\eta\|\bm a\|^2$ stays near one on a silent basis. LIF, ALIF, and LIFRate
units are active ($|a_j|>10^{-9}$) in $0.1\%$ of unit-samples on Full3 and
Hard45, under $0.05\%$ on Roll and Yaw360, and $6$--$20\%$ on Hardstep, so
their gains stay below $0.1\%$ outside Hardstep; CMG units (resting near
$a_j=0.37$ on the logistic tail) and Izhikevich units are active in
${\ge}99.9\%$, and both learn. Persistent activity is thus
necessary for PES here but does not set effort, since Izhikevich--PES shares
it. CMG--PES attains comparable tracking at markedly lower realized effort
and more conservative decoder motion.

% [v10] One two-column figure replaces the two single-column dashboards.
% The 181 mm source was previously scaled into an 88 mm column, a 2.05x
% reduction that left its labels near 3 pt. Drawn at 1:1 here.
\begin{figure*}[t]
    \centering
    \includegraphics[width=\textwidth]{figures/08_Nengo_neuron_family/mechanism_both_platforms.pdf}
    \caption{Mechanism separation across the Nengo neuron-family benchmark,
    X2 (top) and CF2 (bottom). Left: mean $q$-error improvement over the PD
    baseline. Right: normalized residual effort, the RMS residual as a
    fraction of the per-axis limit. Rows are neuron family and learning
    mode, columns are missions; every cell uses the same bounded residual,
    saturation, and allocation interface. Two patterns hold on both
    platforms: PES adaptation, not the neuron model, separates the improving
    rows from the flat ones; and the families that match or exceed
    CMG--PES on error do so at three to twenty times its residual effort.}
    \label{fig:mech_dashboard}
\end{figure*}

%Readout saturation does not imply negligible conductance. In the fresh canonical X2 Full3 rerun, $r_j>0.9$ for $97.7\%$ of unit-samples and mean $c_j\tau$ is $7.08$. Yet $a_j$ clips at $r_j=10^{-5}$.
%Moreover, with $\gamma=2$ and $v_j\ge0$, $r_j\ge1/2$ implies $b_j\ge1/2$ in \eqref{eq:gate_step}; its convex update preserves that high-gate region. Conductance can therefore change membrane voltage without producing a distinguishable decoded activity.

%The resulting decoder has a useful exact reduction. On an interval with
%$\bm a_k=\bar{\bm a}$ and inactive weight projection, let
%$\bm z_k=W_k\bar{\bm a}$. Then \eqref{eq:pes} gives
%\begin{equation}
% \bm z_{k+1}=(1-\lambda_W)\bm z_k+
% \eta\|\bar{\bm a}\|^2[\bm y_k^*-\clip_{[-1,1]}(\bm z_k)].
% \label{eq:constantbasis}
%\end{equation}
%When readout clipping is inactive, this is a first-order target filter.
%For six unit activities its pole is $.995199$, corresponding to about
%$.104$ s. This is a direct-backend identity under the stated conditions,
%not equivalence of the full transient CMG or delayed Nengo dynamics.
%The recorded X2 Full3 activity is identically one after $2.02$ s; the
%reduced recursion matches its saved decoder to $8\times10^{-16}$.

%A focused direct-backend sweep uses $G\in\{1,10^3,10^5,10^6\}$ and $\chi\in\{0,1\}$ on X2 Full3 and CF2 Hardstep. Their reductions remain $56.04$--$56.97\%$ and $18.34$--$18.46\%$, respectively. In contrast, X2 Roll ranges from effectively zero at $G=1$ to $17.36\%$ at $10^5$ and $57.76\%$ at $10^6$ (also tested at $10^2,10^4$). These are fixed-$\eta$ sensitivity tests: before clipping, same-feature PES excitation scales as $\eta G^2\|\bm r\|^2$. They do not isolate neuron expressiveness. Eight selected Nengo checks agree with the direct mean errors to relative difference $1.13\times10^{-7}$.

%\begin{table}[ht]
%\centering\footnotesize
%\caption{Direct-backend controls: mean-error reduction over matched
%PD [\%]. Same $W_0$, PES teacher/update, leak, caps, and allocation;
%canonical CMG uses $G=10^5$. One deterministic run per cell.}
%\label{tab:simplebasis}
%\setlength{\tabcolsep}{4pt}
%\begin{tabular}{lrrr}
%\toprule
%Case & CMG & Paired-linear & Constant\\\midrule
%X2 Roll &17.36&0.73&57.78\\
%X2 Full3 &56.56&6.70&56.98\\
%CF2 Hardstep &18.37&1.50&18.38\\\bottomrule
%\end{tabular}
%\end{table}

\begin{figure}[!t]
\centering
\includegraphics[width=\linewidth]{figures/07_disturbance_and_noise/hover_intermittent.png}
\caption{X2 hover hold under four $0.15$\,N\,m disturbance bursts (shaded)
separated by quiet intervals. Top: attitude error. Centre: residual
magnitude, logarithmic. Bottom: decoder displacement from $W_0$. The
residual is at numerical precision before the first burst. This means
negligible initial correction, not zero computation or energy cost.}
\label{fig:hover_int}
\end{figure}

%%%%%%%%%%%%%%%%%%%%%%%%%%%%%%%%%%%%%%%%%%%%%%%%%%%%%%%%
\subsection{Disturbance Rejection and Retention}
\label{sec:disturbance}
An equilibrium test applies
$\bm\tau_{\rm d}=d_b\rho(t)[\sin(1.31t+.2),\allowbreak
.8\sin(1.73t+1.1),\allowbreak.6\sin(2.17t+2)]\T$ in body axes.
After 5 s of quiet, $\rho$ alternates 5-s on/off intervals with 0.25-s
raised-cosine edge ramps. For X2, $d_b=.15$ N\,m is the first component's
amplitude, not a vector-norm bound.
The four-burst 45-s experiment (Fig.~\ref{fig:hover_int}) cuts burst-phase
mean error from $11.50^\circ$ to $6.65^\circ$ ($42.2\%$) and post-burst
quiet error from $1.21^\circ$ to $.99^\circ$ ($18.0\%$).
The CF2 rerun uses $d_b=.00015$ N\,m (a selected scaling, not an exact
inertia-ratio match); its full-run mean falls from
$.141199$ to $.070259$ rad ($50.24\%$).

%The 180-s X2 record contains eighteen bursts and seventeen complete post-burst quiet intervals, not sixteen bursts. Decoder displacement reaches $.797$ at the final sample, with maximum absolute weight $1.089$.
%One quiet interval is about $4\%$ worse than PD. Continued weight change and nonzero quiet-phase errors do not establish convergence. Because learning and leakage remain active, and the multi-sine phase is not reset between bursts, this study shows persistent adaptation, not causal proof of reuse without relearning. A reset-decoder or frozen-decoder control is needed to test that hypothesis.

\section{Discussion and Conclusion}
\label{sec:conclusion}
% [v7] Rewritten: states the positive result once and plainly, keeps the
% limits in one place instead of after every sentence, and separates the
% analog outlook from the claims.
A bounded residual of six continuous CMG units and eighteen PES-adapted decoder weights augments a fixed quaternion PD baseline without changing its gains or actuator interface, reducing mean quaternion-angle error by $17.3$--$60.6\%$ over ten matched 45-s MuJoCo cases on two modeled platforms, within rotor saturation and with exact baseline recovery at zero authority; direct and Nengo realizations agree to $2.013\times10^{-6}$ relative mean error.

Decoder adaptation is the dominant source of improvement in the tested regime: a fixed decoder gives at most $10.3\%$, and toggling $\chi$ is indistinguishable. Opened gates latch and saturate the readout, where a fixed paired decoder can cancel an axis and PES reduces to lagged proportional action worth about $1.5$--$1.6\times$ the PD stiffness, while sub-threshold units stay in a graded logistic tail; PES shapes the correction on either basis as long as the basis stays active. Across all five neuron families PES is the consistent differentiator, and the dynamics mainly set an effort--accuracy tradeoff. The local certificate gives $\beta_{\max}\sim10^{-4}$--$10^{-3}$, far below the tested $\beta=0.5$.

Limitations: results are in simulation (Sec.~\ref{sec:noise} bounds noise sensitivity; model mismatch and hardware effects are not addressed), the retention horizon is limited, and bounded commands, states, and weights are implementation properties, not global stability proofs.

Future work includes matched-effort comparisons against non-neural bounded bases, less conservative sampled-data certificates, full 6-DOF extension, and, most importantly, validation on physical aerial, ground, and other robotic platforms, to which the same bounded-residual wrapper directly extends.



% Finalize human authors, affiliations, funding, and any required tool-use
% disclosure before submission. No invented author identities are supplied.
\bibliographystyle{IEEEtran}
\bibliography{references}
\end{document}
